\section{Diseño de arquitecturas web escalables}

\subsection{Escalabilidad en aplicaciones web}

La escalabilidad es la capacidad de un sistema para mantener niveles aceptables de rendimiento y disponibilidad ante el incremento de la carga de trabajo, ya sea mediante el aumento de los recursos de un mismo servidor (escalado vertical) o mediante la incorporación de nuevos nodos (escalado horizontal). El uso de balanceo de carga permite distribuir las peticiones entrantes entre múltiples servidores web, maximizando el tiempo de respuesta y el aprovechamiento de recursos, aunque su eficacia depende de evitar la afinidad de servidor y de diseñar componentes sin estado (\textit{stateless}) \cite{guidance_share}. Patrones como el \textit{Intercepting Filter} y el \textit{Front Controller} contribuyen a centralizar el procesamiento de peticiones, facilitando la aplicación de mecanismos transversales de escalabilidad \cite{guidance_share}.

\subsection{Patrones de caché en aplicaciones web}

El almacenamiento en caché constituye una de las estrategias más efectivas para reducir la latencia y mejorar la escalabilidad de una aplicación web, al evitar el recálculo o la recuperación repetida de datos desde el origen. Estudios experimentales han demostrado que el caché en memoria (\textit{in-memory}) puede reducir el tiempo de respuesta hasta en un 62,6\,\% respecto a escenarios sin caché, mientras que el caché basado en archivos y el caché de base de datos ofrecen mejoras de aproximadamente 36,6\,\% y 55,1\,\% respectivamente \cite{researchgate_caching}. La literatura distingue entre el caché a nivel de aplicación —insertado manualmente en el código base y sujeto a errores por su naturaleza transversal— y el caché a nivel de infraestructura o de contenido, evaluándose criterios de reemplazo como \textit{Least Recently Used} (LRU), cuyo rendimiento puede ser superado por estrategias conscientes de la distribución estadística de acceso (Zipf) en más de un 10\,\% de tasa de aciertos \cite{ieee_webcaching}, \cite{acm_caching}. Un estudio cualitativo con diez proyectos de software reales identificó patrones y buenas prácticas para el diseño, implementación y mantenimiento del caché a nivel de aplicación, resaltando su carácter de preocupación transversal (\textit{crosscutting concern}) ajena a la lógica de negocio \cite{ieee_tse_caching}.

\subsection{Documentación y evaluación de arquitecturas web}

La documentación de la arquitectura de software formaliza las decisiones estructurales adoptadas y sus justificaciones, permitiendo su comunicación entre los interesados del proyecto y su evaluación posterior. El método de análisis de arquitecturas basado en atributos de calidad, desarrollado en el ámbito del \textit{Software Engineering Institute} (SEI), propone priorizar los atributos de calidad relevantes para el sistema, diseñar estructuras que reflejen los compromisos (\textit{trade-offs}) entre dichos atributos y, finalmente, evaluar el diseño resultante frente a los requisitos priorizados \cite{arxiv_scientific_sw}. Las \textit{Architecture Decision Records} (ADR) se han consolidado como el formato estándar para documentar estas decisiones, registrando el contexto, la decisión adoptada, las alternativas consideradas y sus consecuencias, lo que evita la repetición de discusiones ya resueltas y facilita la evolución controlada del sistema \cite{tecnovy_arch}.
